% 原书物理页=印刷页33--38；4.2 全节（4.2 末尾的 例1、例2 与 习题 三段实际位于印刷页 39 上部、4.3 节标题之前）.
\section{复表示的特征标}\label{sec:fg-4-2}

从现在起，我们假定群 $G$ 是\emph{有限}的，并且此后只考虑\emph{有限维}的表示.

\begin{definition}\label{def:fg-4-2-1}
设 $(V,\rho)$ 是 $G$ 在一个域 $K$ 上的维数有限的表示. 我们把 $\rho$ 的\newterm{fg-character}{特征标}{character}定义为函数
\[
\chi_\rho=\chi_V:\quad
\begin{array}{rcl}
G & \longrightarrow & K\\
g & \longmapsto & \chi_V(g)=\operatorname{Tr}(\rho(g))
\end{array}
\]
它把 $G$ 的每个元素 $g$ 映到 $V$ 的自同构 $\rho(g)$ 的迹.
\end{definition}

矩阵代数中的经典公式 $\operatorname{Tr}(AB)=\operatorname{Tr}(BA)$ 表明，特征标 $\chi_V$ 满足等式
\[
\forall g\in G,\ \forall h\in G,\ \chi_V(ghg^{-1})=\chi_V(h).
\]
满足这些等式的函数 $G\longrightarrow K$ 称为\newterm{fg-central-function}{中心函数}{central function}，或称\emph{在共轭下不变的函数}.

出于同样的理由，两个同构的表示具有相等的特征标. 显然，正如一个自同态远非由它的迹所确定，我们可以预料，一个表示的特征标是它的一个相当粗糙的不变量. 然而我们将看到，在域 $K=\mathbb{C}$ 上，具有相同特征标的两个表示是同构的！在此之前，我们注意到特征标知道表示的维数，因为 $\chi_V(e)=\dim_K(V)$.

\begin{remark}\label{rem:fg-4-2-2}
在文献中，``$G$ 的特征标''这个名称常被用来指从 $G$ 到乘法群 $K^\times$ 的同态. 这种意义下的一个特征标是上面意义下的维数 $1$ 的表示的特征标. 但维数 $>1$ 的表示 $V$ 的特征标 $\chi_V$ 决不是乘性的（由于它在 $e$ 处取值 $\dim(V)$）. 因此这里有术语上的冲突.
\end{remark}

\begin{theorem}\label{thm:fg-4-2-3}
设 $V$ 与 $W$ 是 $G$ 的两个维数有限的表示. 那么：
\begin{enumerate}[label=\roman*)]
\item 逆步表示 $V^*$ 的特征标由 $\chi_{V^*}(g)=\chi_V(g^{-1})$ 给出；
\item 表示 $V\oplus W$ 的特征标由 $\chi_{V\oplus W}(g)=\chi_V(g)+\chi_W(g)$ 给出；
\item 表示 $V\otimes W$ 的特征标由 $\chi_{V\otimes W}(g)=\chi_V(g)\chi_W(g)$ 给出；
\item 表示 $H=\operatorname{Hom}(V,W)$ 的特征标由 $\chi_H(g)=\chi_W(g)\chi_V(g^{-1})$ 给出.
\end{enumerate}
\end{theorem}

\begin{proof}
设 $B_V=\{e_i,\ i\in I\}$ 是 $K$-向量空间 $V$ 的一组基，而 $B_{V^*}=\{e_i^*,\ i\in I\}$ 是它在 $V^*$ 中的对偶基. 由迹的定义，有
\[
\chi_V(g)=\sum_{i\in I}\langle ge_i,\,e_i^*\rangle.
\]
由此可得
\[
\chi_V(g)=\sum_{i\in I}\langle ge_i,\,e_i^*\rangle=\sum_{i\in I}\langle e_i,\,g^{-1}e_i^*\rangle=\chi_{V^*}(g^{-1}),
\]
其中中间的等式来自逆步表示的定义. 这就证明了 i).

再取定 $W$ 的一组基 $B_W=\{\varepsilon_j,\ j\in J\}$. 并集 $B_V\cup B_W$ 是 $V\oplus W$ 的一组基，而对偶基的并集 $B_{V^*}\cup B_{W^*}$ 是它在 $(V\oplus W)^*=V^*\oplus W^*$ 中的对偶基. 于是，按直和表示的定义，有
\[
\chi_{V\oplus W}(g)=\sum_{i\in I}\langle ge_i,\,e_i^*\rangle+\sum_{j\in J}\langle g\varepsilon_j,\,\varepsilon_j^*\rangle=\chi_V(g)+\chi_W(g).
\]
另一方面，前面已经看到集合 $B_{V\otimes W}:=\{e_i\otimes\varepsilon_j,\ i\in I,\ j\in J\}$ 是 $V\otimes W$ 的一组基. 存在一个明显的 $K$-向量空间同构 $V^*\otimes W^*\xrightarrow{\ \sim\ }(V\otimes W)^*$，它把张量 $v^*\otimes w^*$ 映到 $V\otimes_K W$ 上与双线性形式 $(v,w)\mapsto\langle v,v^*\rangle\langle w,w^*\rangle$ 相伴的那个线性形式. 这使我们能把集合 $\{e_i^*\otimes\varepsilon_j^*,\ i\in I,\ j\in J\}$ 等同于 $B_{V\otimes W}$ 的对偶基. 于是有
\[
\begin{aligned}
\chi_{V\otimes W}(g) &= \sum_{i,j}\langle ge_i\otimes g\varepsilon_j,\ e_i^*\otimes\varepsilon_j^*\rangle=\sum_{i,j}\langle ge_i,\,e_i^*\rangle\langle g\varepsilon_j,\,\varepsilon_j^*\rangle\\
&=\Biggl(\sum_i\langle ge_i,\,e_i^*\rangle\Biggr)\Biggl(\sum_j\langle g\varepsilon_j,\,\varepsilon_j^*\rangle\Biggr)=\chi_V(g)\chi_W(g).
\end{aligned}
\]
最后，前面已经看到 $\operatorname{Hom}(V,W)\simeq V^\times\otimes W$\footnote{译注：原文此处把逆步表示记作 $V^\times$，即前文记号下的 $V^*$（参见 \ref{np:fg-4-1-21}--\ref{lemma:fg-4-1-22} 段，那里得到 $\operatorname{Hom}_K(V,W)$ 上的表示同构于 $W\otimes V^*$）.}，故 iv) 由 i) 与 iii) 推出.
\end{proof}

\begin{example}\label{ex:fg-4-2-4}
设 $X$ 是一个 $G$-集，$KX$ 是与之相伴的置换表示. 它的特征标由下式给出：
\[
\chi_{KX}(g)=\ \text{$g$ 在 $X$ 中不动点的个数}.
\]
\end{example}

\begin{example}\label{ex:fg-4-2-5}
设 $G=D_n$，$V=\mathbb{C}^2$ 是标准表示. 那么，若 $g$ 不属于 $D_n$ 中由转角为 $\frac{2\pi}{n}$ 的旋转 $r_n$ 生成的循环正规子群 $C_n$，则 $\chi_V(g)=0$；而对 $k=0,\cdots,n-1$，有 $\chi_V(r_n^k)=2\cos(2k\pi/n)$.
\end{example}

\begin{numpara}[复的情形]\label{np:fg-4-2-6}
从此我们假定 $K=\mathbb{C}$. 我们用 $z^*$ 表示复数 $z$ 的共轭.

$\mathbb{C}$ 的第一个特殊性在于它是代数闭的. 于是自同构 $\rho(g)$ 的迹是它的\emph{特征值的和}. 由于 $\rho(g)^{|G|}=\operatorname{Id}_V$，$\rho(g)$ 的特征值都是单位根，因而有 $\chi_V(g^{-1})=\chi_V(g)^*$.

$\mathbb{C}$ 的第二个特殊性来自 $zz^*$ 对一切 $z$ 的正性. 这使得我们能在 $G$ 上复函数组成的空间 $\mathbb{C}G=\mathbb{C}[G]=\mathcal{F}(G,\mathbb{C})$ 上赋予其 $L^2$ 范数，即 Hermite 内积：
\[
(f,f')=|G|^{-1}\sum_{g\in G}f(g)f'(g)^*.
\]
我们也将考虑它的子空间
\[
\mathcal{F}^{cent}(G,\mathbb{C})=\{f\in\mathcal{F}(G,\mathbb{C}),\ \forall g,h\in G,\ f(hgh^{-1})=f(g)\},
\]
即中心函数组成的空间，并赋予其诱导范数. （由于这里的 $G$ 是有限的，我们已从记号 $\mathcal{F}_0$ 中略去了下标 0.）下面的\newterm{fg-orthogonality-relations}{正交关系}{orthogonality relations}是特征标理论的钥匙.
\end{numpara}

\begin{theorem}\label{thm:fg-4-2-7}
设 $V$ 与 $W$ 是 $G$ 的两个不可约表示. 那么：
\begin{enumerate}[label=\roman*)]
\item 若 $W$ 不同构于 $V$，则 $(\chi_V,\chi_W)=0$；
\item 若 $W$ 同构于 $V$，则 $(\chi_V,\chi_V)=1$.
\end{enumerate}
\end{theorem}

\begin{proof}
先考虑任意一个表示 $(\rho,V)$，并考察 $V$ 在 $G$ 作用下由固定向量组成的子空间 $V^G$. $V$ 的自同态 $p=|G|^{-1}\sum_{g\in G}\rho(g)$ 是 $V$ 到 $V^G$ 上的一个投影算子. 于是它的迹为
\[
\operatorname{Tr}(p)=\dim(V^G)=|G|^{-1}\sum_{g\in G}\operatorname{Tr}(\rho(g))=|G|^{-1}\sum_{g\in G}\chi_V(g)=(\chi_V,\mathbf{1}),
\]
这里 $\mathbf{1}$ 表示 $G$ 上取值恒为 1 的常值函数.

回到定理陈述中的记号，并把这个结果应用到表示 $H=\operatorname{Hom}(V,W)$ 上. 我们有 $(\chi_H,\mathbf{1})=(\chi_W\chi_V^*,\mathbf{1})=(\chi_W,\chi_V)$. 根据 Schur 引理（命题~\ref{prop:fg-4-1-8}），$K$-向量空间 $\operatorname{Hom}(V,W)^G$ 的维数在 $V$ 与 $W$ 不同构（相应地，同构）时为 0（相应地，1）. 于是 $(\chi_H,\mathbf{1})=0$ 或 $1$，视我们处于情形 i) 还是情形 ii) 而定，正如所述.
\end{proof}

\begin{corollary}\label{cor:fg-4-2-8}
设 $V$ 是 $G$ 的一个表示，且 $V_1^{\oplus a_1}\oplus\ldots\oplus V_k^{\oplus a_k}$ 是它的同型分解.
\begin{enumerate}[label=\roman*)]
\item 对一切 $i=1,\ldots,k$，$V_i$ 在 $V$ 中的重数 $a_i$ 由 $a_i=(\chi_V,\chi_{V_i})$ 给出. 特别地，两个表示同构当且仅当它们有相同的特征标.
\item 我们有 $(\chi_V,\chi_V)=\sum_{i=1}^k a_i^2$. 特别地，$V$ 不可约当且仅当 $(\chi_V,\chi_V)=1$.
\item $G$ 的任何不可约表示 $W$ 都出现在 $G$ 的正则表示中，其重数等于它的维数 $\dim(W)$. 特别地，$|G|=\sum_W\dim(W)^2$，其中 $W$ 取遍 $G$ 的不可约表示的同构类集合.
\end{enumerate}
\end{corollary}

\begin{proof}
i)$+$ii) 由正交关系以及 $\chi_{V_i^{\oplus a_i}}=a_i\chi_{V_i}$ 这一事实推出.
iii) 正则表示 $R$ 的特征标 $r$ 满足 $r(1)=|G|$，且对一切 $g\neq1$ 有 $r(g)=0$. 于是对 $G$ 的任何表示 $W$，有
\[
(r,\chi_W)=|G|^{-1}(r(1)\chi_W(1)^*+0)=\dim(W).
\]
设 $W$ 不可约，且 $R$ 同构于 $W_1^{\oplus a_1}\oplus\ldots\oplus W_k^{\oplus a_k}$；那么由正交关系即可推出 $W$ 同构于某个 $W_i$，且 $a_i=\dim(W)$.
\end{proof}

最后这个结论表明，$G$ 的不可约表示的同构类只有有限多个. 事实上，恰好有 $cl(G)$ 个（即 $\mathcal{F}^{cent}(G,\mathbb{C})$ 的维数）. 理由如下：

\begin{theorem}\label{thm:fg-4-2-9}
$G$ 的互不同构的不可约表示的特征标构成 $G$ 上中心函数空间的一组基. 因此它们的个数等于群 $G$ 的共轭类的个数.
\end{theorem}

\begin{proof}
要证的是：使得对一切不可约表示 $V$ 都有 $(f,\chi_V)=0$ 的中心函数 $f$ 只有零函数. 对这样一个 $V$，令 $\varphi_V=\sum_{g\in G}f(g)\rho_V(g)$. 由于 $f$ 是中心函数，可以看出 $\varphi$ 是 $V$ 到 $V$ 的一个 $G$-态射，故由 Schur 引理它是一个位似；又由于其迹 $|G|(f,\chi_V^*)$ 为零，便有 $\varphi_V=0$. 于是正则表示的空间 $\mathcal{F}(G,\mathbb{C})$ 上的自同态 $\sum_{g\in G}f(g)R(g)$ 恒等于零. 由于它的向量 $R(g)(e_1)=e_g$ 是 $\mathbb{C}$-线性无关的，所以 $f(g)$ 对一切 $g\in G$ 确实为零.
\end{proof}

\begin{remark}\label{rem:fg-4-2-10}
不可约表示的特征标构成中心函数空间的一组基这一事实，在系数域为某个域 $K$ 时仍然成立，只要 $|G|$ 在 $K$ 中可逆，并且 $K$ 包含一个 $|G|$ 次本原单位根. 否则便有反例. 最引人注目的情形是 $K$ 的特征为 $p$ 而 $G$ 是一个 $p$-群的情形. 此时任何表示的特征标都是常值！

一个表示由其特征标确定（在同构意义下）这一事实在特征零的域上仍然成立. 但当 $K$ 的特征为正的 $\ell$ 时它总是不成立的. 例如维数为 $\ell$ 的平凡表示有零特征标！同样，特征标也看不出 $\rho\oplus\sigma$ 与 $\rho$ 被 $\sigma$ 的一个非分裂扩张之间的差别.
\end{remark}

\begin{numpara}[特征标表]\label{np:fg-4-2-11}
根据前文，我们可以给一个有限群 $G$ 联系上一张\newterm{fg-character-table}{特征标表}{character table}，即一张复数表，其各行以 $G$ 的不可约表示 $\rho$ 的同构类为指标，各列以 $G$ 中的共轭类 $c$ 为指标. 在行 $\rho$ 与列 $c$ 处，我们放上 $\chi_\rho(x)$，其中 $x$ 是 $c$ 中的任意一个元素. 下面是几个例子.

\emph{循环群 $G=\mathbb{Z}/n\mathbb{Z}$.} 由于这个群是交换群，它的共轭类都是单元素集，即简单地就是 $\{0\},\cdots,\{n-1\}$. 另一方面，前面已经看到 $G$ 的不可约表示都是 1 维的，由同态 $\rho_{\bar{\imath}}:\ \bar{j}\in G\mapsto\omega^{\bar{j}\bar{\imath}}$ 给出，其中 $\omega$ 是取定的一个 $n$ 次本原单位根，而 $\bar{\imath}$ 取遍模 $n$ 的整数. 于是这张表简单地由公式 $\chi_{\rho_{\bar{\imath}}}(\bar{j})=\omega^{\bar{\imath}\bar{j}}$ 给出. 尽管如此要注意，这张表的书写方式依赖于 $\omega$ 的选取；改变 $\omega$ 将诱导表的各行之间的一个置换.

\emph{群 $S_3$.} 共轭类是 $\{e\}$、$c_2=\{(1,2),(2,3),(1,3)\}$ 与 $c_3=\{(1,2,3),(2,1,3)\}$. 我们需要找出 3 个不可约表示. 我们已经知道平凡的 $\mathbf{1}$，以及符号特征标 $\operatorname{Sgn}:S_3\longrightarrow\{\pm1\}$. 为了找到第三个，可以想到 $S_3=D_3$ 是保持一个正三角形不变的等距所成的群，于是有维数 2 的标准表示 $\mathrm{Std}$：$S_3=D_3\subset\mathrm{O}_2(\mathbb{R})\subset\mathrm{GL}_2(\mathbb{R})\subset\mathrm{GL}_2(\mathbb{C})$. 现在可以列出表：
\[
\begin{array}{c|c|c|c}
 & \{e\} & c_2 & c_3\\\hline
\mathbf{1} & 1 & 1 & 1\\\hline
\operatorname{Sgn} & 1 & -1 & 1\\\hline
\mathrm{Std} & 2 & 0 & -1
\end{array}
\]
\emph{群 $D_4$.} 它形如 $\langle r\rangle\rtimes\langle s\rangle$，其中 $r$ 是转角为 $\pi/2$ 的旋转（因而生成一个同构于 $\mathbb{Z}/4\mathbb{Z}$ 的子群），$s$ 是以 $(Ox)$ 为轴的正交反射，二者都保持 $\mathbb{R}^2$ 中以 $(\pm1,\pm1)$ 为顶点的正方形. 从关系 $srs=r^2$\footnote{译注：原文如此；按 $D_4$ 的通常定义，此关系似应为 $srs=r^{-1}=r^3$.} 出发，可以找到共轭类 $\{e\}$、$\{r^2\}$、$\{r,r^3s\}$\footnote{译注：原文如此；此集似应为 $\{r,r^3\}$.}、$\{s,r^2s\}$ 与 $\{rs,r^3s\}$. 为了找出不可约表示，注意到 $D_4/\langle r^2\rangle$ 同构于 $\mathbb{Z}/2\mathbb{Z}\rtimes\mathbb{Z}/2\mathbb{Z}=\mathbb{Z}/2\mathbb{Z}\times\mathbb{Z}/2\mathbb{Z}$. 由此得到 4 个维数 1 的特征标 $\mathbf{1}$、$\chi$、$\psi$ 与 $\chi\psi$，其定义为 $\chi(r)=-1$、$\chi(s)=1$，以及 $\psi(r)=1$、$\psi(s)=-1$. 只剩一个表示要找，那就是维数 2 的标准表示. 现在可以写下表：
\[
\begin{array}{c|c|c|c|c|c}
 & \{e\} & \{r^2\} & \{r,r^3s\} & \{s,r^2s\} & \{rs,r^3s\}\\\hline
\mathbf{1} & 1 & 1 & 1 & 1 & 1\\\hline
\chi & 1 & 1 & -1 & 1 & -1\\\hline
\psi & 1 & 1 & 1 & -1 & -1\\\hline
\chi\psi & 1 & 1 & -1 & -1 & 1\\\hline
\mathrm{Std} & 2 & -2 & 0 & 0 & 0
\end{array}
\]
\emph{群 $H_8$.} 容易验证共轭类是 $\{1\}$、$\{-1\}$、$\{\pm i\}$、$\{\pm j\}$ 与 $\{\pm k\}$. 另一方面，$H_8/\{\pm1\}$ 同构于 $\mathbb{Z}/2\mathbb{Z}\times\mathbb{Z}/2\mathbb{Z}$，其生成元可取 $i$ 与 $j$ 的像（例如）. 由此得到 4 个维数 1 的特征标 $\mathbf{1}$、$\chi$、$\psi$ 与 $\chi\psi$，其定义为 $\chi(i)=-1$、$\chi(j)=1$，以及 $\psi(i)=1$、$\psi(j)=-1$. 剩下还要找出最后一个不可约表示. 考虑 $\mathbb{H}=\mathbb{C}\oplus\mathbb{C}j$，把它看作一个维数 2、以 $1,j$ 为基的 $\mathbb{C}$-向量空间，并让 $H_8$ 通过左乘作用在 $\mathbb{H}$ 上：$\rho(g)(x):=xg^{-1}$. 这样便得到一个 $\mathbb{C}$-线性表示（注意右乘将不是 $\mathbb{C}$-线性的，而只是 $\mathbb{R}$-线性的）\footnote{译注：原文此句中``左乘/右乘''二字疑互换：公式 $\rho(g)(x)=xg^{-1}$ 实为右乘，且由 $\mathbb{H}$ 的结合律，右乘总是 $\mathbb{C}$-线性的；反而是左乘一般只是 $\mathbb{R}$-线性的（例如左乘 $j$ 不与左乘 $i$ 交换，因 $ji=-k\ne k=ij$）.}. 于是可以列出下表：
\[
\begin{array}{c|c|c|c|c|c}
 & \{1\} & \{-1\} & \{\pm i\} & \{\pm j\} & \{\pm k\}\\\hline
\mathbf{1} & 1 & 1 & 1 & 1 & 1\\\hline
\chi & 1 & 1 & -1 & 1 & -1\\\hline
\psi & 1 & 1 & 1 & -1 & -1\\\hline
\chi\psi & 1 & 1 & -1 & -1 & 1\\\hline
\mathbb{H} & 2 & -2 & 0 & 0 & 0
\end{array}
\]
\emph{注 1：}如果找不到最后一个不可约表示，总可以从推论~\ref{cor:fg-4-2-8}~iii) 推出特征标表的最后一行. 因为它说的正是：第一列的欧几里得范数的平方等于群 $G$ 的阶，而其余各列都与第一列正交. 反过来，这个关系可以用来发现一张表中可能存在的错误.

\emph{注 2：}前面两个例子表明，两个不同构的群 $G,G'$ 可以具有``同构的''特征标表，意思是：存在从 $G$ 的共轭类到 $G'$ 的共轭类（相应地，从 $G$ 的不可约表示类到 $G'$ 的不可约表示类）的一个双射，使得所得的表相同. 不过，仍然存在两种方式，仅从表出发就能区分两个具有同一张表的群，\emph{只需记得一些结构}.

\emph{例 1：}映射 $g\mapsto g^2$ 与共轭相容，因而诱导出从 $G$ 的共轭类集合到其自身的一个映射 $c\mapsto c^2$. 对于 $D_4$，我们注意到 $\{s,r^2s\}^2=\{e\}$，而对于 $H_8$，则有 $\{\pm j\}^2=\{-1\}$. 这说明，使特征标表相同的那种共轭类之间的双射并不与``共轭类的平方''这一运算相容，这就使我们能够区分这两个群.

\emph{例 2：}与其用共轭类上的一个运算（如上面的``平方''映射），不如用表示上的一个运算. 设 $(V,\rho)$ 是 $G$ 的一个表示，我们把它的\newterm{fg-symmetric-square}{对称平方}{symmetric square}记作
\[
\operatorname{Sym}^2V:=(V\otimes V)\big/\operatorname{Vect}\{(v_1\otimes v_2-v_2\otimes v_1),\ v_1,v_2\in V\}.
\]
\emph{习题：}验证这是 $\rho\otimes\rho$ 的一个商表示，记作 $\operatorname{Sym}^2(\rho)$，并验证其特征标由 $\chi_{\operatorname{Sym}^2(\rho)}(g)=\frac{1}{2}\bigl(\chi_\rho(g)^2+\chi_\rho(g^2)\bigr)$ 给出. 由此推出 $\operatorname{Sym}^2(\mathrm{Std})=\mathbf{1}+\chi+\chi\psi$，而 $\operatorname{Sym}^2(\mathbb{H})=\chi+\psi+\chi\psi$. 这说明 $D_4$ 的特征标与 $H_8$ 的特征标之间的那个双射并不与``对称平方''这一运算相容，这就使我们能够区分这两个群.
\end{numpara}
